\section*{Problem 9}

A trusted entity chooses primes \(p,q\), sets \(n=pq\), and chooses integers \(k_1,k_2,k_3\)
satisfying
\[
k_1 k_2 k_3 \equiv 1 \pmod{(p-1)(q-1)}.
\]
Keys:
\[
\text{Alice: }(n,k_1,k_2),\qquad
\text{Bob: }(n,k_2,k_3),\qquad
\text{Carla: }(n,k_1,k_3).
\]

\subsection*{(a)}
Alice wants Bob and Carla to decrypt the same message \(m\).
She encrypts using the key component they \emph{both} possess:
Bob has \(k_2\), Carla also has \(k_2\)?
No—Carla has \(k_1,k_3\).  
Bob has \(k_2,k_3\).

The common exponent between Bob and Carla is \(k_3\).

Thus Alice sends:
\[
c \equiv m^{k_3} \pmod{n}.
\]

Bob decrypts using \(k_2\) since \(k_2 k_3 k_1 \equiv 1\).  
Carla decrypts using \(k_1\).

\subsection*{(b)}
Alice wants only Carla to decrypt. She should use an exponent that:
- Carla knows,
- Bob does \emph{not} know.

Carla knows \(k_1\) and \(k_3\); Bob knows \(k_2\) and \(k_3\).  
So the exponent Bob does not know is \(k_1\).

Alice sends:
\[
c \equiv m^{k_1} \pmod{n}.
\]
Carla decrypts using \(k_2 k_3\), but Bob cannot decrypt.
